\section{Introduction}
\label{sec:intro}

Language models are starting to operate as security agents rather than only as
advisers. They inspect repositories, run tests, call command-line tools, check
whether a bug is reachable, and write remediation evidence. That shift makes
safety alignment harder to reason about. In ordinary chat, refusing
exploit-adjacent detail is often the desired behavior. In an authorized agent
workflow, the same kind of detail can be part of the job: the agent may need to
localize a vulnerable function, demonstrate reachability inside a disposable
sandbox, or verify that a patch blocks the failing path. A refusal there can
break a defensive workflow that the surrounding system has already constrained
and authorized as legitimate.

\looseness=-1
This paper evaluates that tension empirically. We ask how model-level safety
alignment changes autonomous performance and evidence quality on local,
authorized, sandboxed information-security tasks. The question is not whether
safety controls are necessary; they are. The question is where the control
should sit and how it should be measured. Model-level refusal can reduce
open-ended misuse, but be too coarse for contained workflows requiring local
construction and validation. Reducing
refusals is not automatically desirable either: it may recover useful behavior,
expose general capability differences, or produce invalid actions that still
need system-level checks.

\looseness=-1
We build one paired local evaluation around four aligned versus
less-restricted model pairs. Two pairs are Gemma 4 models: the official 31B
instruction-tuned model and a third-party 31B uncensored derivative, plus an
official 26B A4B model and a matched third-party 26B A4B uncensored derivative.
The same evaluation includes Qwen2.5-Coder 7B and Llama 3.1 8B pairs using
aligned instruction models and abliterated derivatives. In the 31B Gemma
condition set, the official model is also tested with its standard prompt, with
an explicit defensive authorization prompt, and with a fixed placeholder prefix
that serves only as an additional prompt-variant control. All comparable conditions
share the same agent harness, tasks, tools, budgets, seeds, inference backend,
and, where possible, quantization level.

\looseness=-1
The security evaluation contains 1,500 saved runs across the four model pairs.
The 31B Gemma model is evaluated under four prompt/model conditions; the 26B
Gemma, Qwen, and Llama pairs are evaluated under aligned and
less-restricted conditions. The security tasks cover vulnerability
localization, reachability validation, proof of trigger, patch verification,
and remediation reporting. Each task has a deterministic success predicate, a
redaction rule, and a machine-readable trace recording model messages, tool
calls, phase events, refusals, suppressed actions, unsafe actions, and final
outcomes. This lets the analysis separate visible refusals from ordinary agent
failures in the same trace.

\looseness=-1
Raw success rates are not enough. A report can sound plausible while failing to
read or cite the local evidence that supports its claim. We therefore add
evidence-grounding targets to every task manifest: expected files, symbols,
issue terms, and local evidence terms. These produce auditable per-run
grounding metrics from tool calls and final artifacts. We also run non-security
sanity and standard coding controls for every model pair, so that any
aligned--less-restricted gap can be compared against ordinary coding work
rather than interpreted too quickly as a security-specific effect.

\looseness=-1
The cross-family result is deliberately less tidy than a single-family account.
Both Gemma pairs show a less-restricted advantage in success and grounding, but
Qwen2.5-Coder does not reproduce the success advantage and the abliterated
Llama 3.1 derivative performs much worse under the tool protocol. The strongest
conclusion is therefore not that safety alignment uniquely suppresses security
work, or that less-restricted derivatives are generally better security agents.
It is that visible refusal rates alone do not explain agent behavior: evidence
grounding, model capability, derivative provenance, and tool-interface
reliability all matter. No family passes the hard proof and patch-verification
checks, partly because proof-of-trigger checks require an unprompted term
(Section~\ref{sec:results}).

The paper makes three contributions.
\begin{enumerate}[leftmargin=*]
  \item \looseness=-1
  We introduce a system-level benchmark for authorized security agents:
  30 local sandboxed tasks with fixed tools, deterministic success checks,
  redaction rules, and trace labels that distinguish refusal from tool,
  localization, validation, and reporting failures.\footnote{Code and data:
  \codeurl.}
  \item \looseness=-1
  We run a single cross-family evaluation across Gemma 4 31B, Gemma 4
  26B A4B, Qwen2.5-Coder 7B, and Llama 3.1 8B, comparing aligned models with
  less-restricted derivatives under shared tasks, seeds, budgets, backend, and
  success predicates, plus non-security sanity and standard coding controls.
  \item \looseness=-1
  We show that the observed behavior is not explained by visible refusal
  alone. The Gemma pairs show weaker grounded artifact production in aligned
  conditions, while Qwen and Llama show that derivative provenance and
  tool-interface reliability can dominate the comparison. Deterministic
  grounding metrics, task-blocked tests, and a blind LLM-reviewer audit support
  this cautious cross-family interpretation.
\end{enumerate}

\looseness=-1
The scope is intentionally limited. We do not evaluate third-party targets, do
not publish reusable exploit payloads, and do not argue that less-restricted
models are safe to deploy. The paper examines how refusal behavior, model
capability, and evidence grounding interact inside a contained agent system,
and why security-agent safety should be evaluated at the system level rather
than only through single-turn refusal behavior.
